\section{Related Work}
\label{sec:related}

A recent survey~\cite{harlow2024survey} catalogues the rapidly growing use
of mmWave radar in robotics across motion estimation, localization, and
mapping; we focus on the radar-inertial odometry thread most relevant to
single-chip aerial fusion.

\textbf{Radar ego-velocity estimation.}
Doppler-based ego-velocity estimation from mmWave radar was
demonstrated by Kellner et al.~\cite{kellner2014instantaneous} and
extended to 3D with weighted least squares by
Kramer et al.~\cite{kramer2020rve}, whose RVE system fuses radar
with \ac{imu} in a discrete-time \ac{ekf}.  Doer and
Trommer~\cite{doer2020ekf} present an \ac{ekf}-based \ac{rio} with online
extrinsic calibration. Our work uses the same per-frame \ac{wls}
ego-velocity estimator but embeds it in a continuous-time
batch/sliding-window optimizer.

\textbf{Continuous-time trajectory estimation.}
Furgale et al.~\cite{furgale2013unified} introduced uniform B-splines
for continuous-time visual-inertial calibration.  The basalt
\ac{vio}~\cite{usenko2020basalt} uses cumulative B-splines on Lie groups
for stereo-inertial odometry; we adapt this basalt-style
cumulative SO(3) B-spline to radar-inertial fusion.

\textbf{Marginalization in sliding-window estimators.}
OKVIS~\cite{leutenegger2015keyframe} and
VINS-Mono~\cite{qin2018vins} pioneered Schur complement
marginalization for visual-inertial systems, propagating information
from dropped keyframes as a dense Gaussian prior. We apply the same
principle to continuous-time B-spline control points and document
the practical challenges of prior scaling when the marginal's
orientation and position information differ by ten orders of magnitude.

\textbf{Closest related work.}
Continuous-time radar-inertial odometry was introduced by
Ng et al.~\cite{ng2021ctrio}, who fuse Doppler measurements from
\emph{multiple automotive} radars with an \ac{imu} as a least-squares
problem over a B-spline trajectory, evaluated on passenger vehicles.
Burnett et al.~\cite{burnett2024gp} use a Gaussian-process motion
prior for spinning-radar- and lidar-inertial odometry, and
\v{S}tironja et al.~\cite{stironja2026rio} model the \ac{imu} trajectory
continuously inside an \ac{rio} estimator with online spatio-temporal
calibration.  Single-chip radar-inertial estimation has itself
received recent attention: Huang et al.~\cite{huang2024lessismore}
exploit radar cross-section and Doppler physics to make a sparse
single-chip cloud usable in a discrete-time sliding-window
optimization, evaluated on ground-robot and handheld platforms.
Our setting differs from all of these in sensor and
regime: a \emph{single-chip} radar (10--60 points per frame,
2-TX elevation array) on a quadrotor in aggressive racing and
attitude-tumbling flight up to \SI{10}{\radian\per\second}, a
regime in which, as \autoref{sec:backflips} shows, the radar
measurement model itself degrades in ways invisible on ground
vehicles and in benign aerial navigation.

\textbf{Single-chip \ac{rio} on aerial platforms.}
Single-chip radar-inertial fusion has been shown on multirotors, but only
for benign flight.  Michalczyk et al.~\cite{michalczyk2022ekf} use a
discrete-time \ac{ekf} with scan-to-scan point matching and Doppler (a
factor-graph reformulation follows in~\cite{michalczyk2024fg}); Girod et
al.~\cite{girod2024brio} run a real-time onboard M-estimator smoother on a
quadrotor in cities and forests, but cap velocity at
\SI{2}{\metre\per\second} and anchor the vertical channel with a
barometer.  None uses a continuous-time spline, addresses aggressive or
tumbling flight, or, as we do, recovers the vertical channel from radar
and \ac{imu} alone (\autoref{sec:baselines}).

\textbf{Sliding-window marginalization and consistency.}
That na\"ive sliding-window marginalization is inconsistent is long
established in discrete time~\cite{dongsi2012consistency}, where \ac{fej}
methods fix linearization points to preserve the observability nullspace.
In continuous time, CLIC~\cite{lv2023clic} and
LIO-MARS~\cite{quenzel2025liomars} marginalize spline knots leaving the
window by adjacency to the dropped knots; this is correct when the
estimated states are all local to the window, but not when a
\emph{global} shared \ac{imu} bias couples every \ac{imu} factor and
tempts the adjacency rule to pull in far-edge factors, producing an
over-confident prior.  Our contribution, narrower than and complementary
to \ac{fej}, is a factor-\emph{set} selection rule that exploits B-spline
local support to keep the marginal exactly closed despite this shared
bias, together with a diagnostic symptom pattern for the inconsistent
variant; both are developed in \autoref{sec:sw} and
\autoref{sec:prior_scale}.
In the discrete-time UAV line, the \ac{ekf}-\ac{rio} family of Doer and
Trommer~\cite{doer2020ekf} defines the state of practice.
Absolute \ac{rmse} is not comparable across papers: our own quantization
ablation (\autoref{sec:heldout}) shows firmware configuration alone
changes Doppler resolution by $12\times$, so rather than quote their
numbers we cross-validate against this line directly:
\autoref{sec:baselines} runs our system on the
public ICINS-2021 flight datasets of~\cite{doer2021yrio} and re-evaluates
their \ac{ekf} estimators under our causal metrics on the same data, and
reports the regime-dependent result of porting their filters to our
platform.
